\scanpage{154}
\[
\deleted{\widetilde\rho^{3}(\widetilde\tau'_\infty)
=\widetilde\omega_0^{-1}\widetilde\tau'_\infty\widetilde\omega_0
=\widetilde\omega_0^{-1}\underbrace{\widetilde\tau'_\infty\widetilde\omega_0\widetilde\tau_\infty^{\prime-1}}\widetilde\tau'_\infty
=\widetilde\omega_0^{-2}\widetilde\tau'_\infty
=\ell_0^{\prime-1}\widetilde\tau'_\infty}
\]
\[
\tau'',\ \tau,\ \unsure{\varepsilon|\xi}_0,\ \unsure{\varepsilon|\xi}_1,\ \unsure{\varepsilon|\xi}_2
\]
\diagram{\[\begin{array}{lll}
\tau'',\unsure{\varepsilon|\xi}_0,\unsure{\varepsilon|\xi}_1 & \deleted{\unsure{engendrent}}\ \text{engendrent} & Gl(2,\mathbb Z)\Bigg]\ \tau\ \text{permute}\\
\tau'',\unsure{\varepsilon|\xi}_2,\unsure{\varepsilon|\xi}_1^{-1} & \rule{4cm}{.4pt} & Gl(2,\mathbb Z)\qquad\text{ces deux ss-}\\
\tau'',\unsure{\varepsilon|\xi}_1 & \rule{4cm}{.4pt} & D_{\deleted{\infty}\unsure{\varepsilon|\xi}_1}\qquad\text{groupes }Gl(2,\mathbb Z)
\end{array}\]}
\margin{et opère $-1$ sur $D_{\unsure{\varepsilon|\xi}_1}$}
\[
\mathcal T^!_{1,2}\ \overset{?}{\overset{\sim}{\longleftarrow}}\ Gl(2,\mathbb Z)*_{D_\infty}Gl(2,\mathbb Z)
\]

\figure{\begin{tikzpicture}[x=1cm,y=1cm]
\draw (-2.6,-.2) .. controls (-2.8,1.5) and (-1.8,3.1) .. (-.5,3.35) .. controls (1.5,4.05) and (2.7,2.1) .. (2.55,-.5);
\draw (-2.65,-.35) .. controls (-2.6,-1.65) and (-1.5,-2.65) .. (-.2,-2.75) .. controls (.8,-2.9) and (1.25,-2.55) .. (1.35,-2.4);
\draw (1.48,-2.28) .. controls (2.05,-1.95) and (2.38,-1.4) .. (2.45,-.85);
\draw (-.16,3.4) .. controls (-.7,2.65) and (-.88,1.1) .. (-.99,.08) -- (-.94,.02) -- (-.93,.2);
\draw (.33,3.28) .. controls (-.1,2.55) and (-.34,1.3) .. (-.37,.06);
\draw (.58,3.16) .. controls (.8,2.95) and (.78,1.3) .. (.97,.36);
\draw (.93,3.34) .. controls (1.2,2.45) and (1.45,1.15) .. (1.55,.42);
\draw (-2.32,1.85) -- (-1.8,2.48) -- (-1.8,2.35);
\end{tikzpicture}}

$3g$ cylindres qui sont deux à deux en posit°, standart (deux cylindres polaires $\equiv$ pôle, ou un cylindre polaire et un transverse relatif $\unsure{\leftrightarrow|\longrightarrow}\ \equiv$ axe équatorial) ou ne se rencontrent pas.
\[
2g+2\,\frac{g(g-1)}2=2g\left(1+\frac{g-1}2\right. =2\,\frac{g(g+1)}2
\]
paires de cylindres en position ``standart'', et les autres paires sont disjointes

MR LE PROFESSEUR GROTHENDIECK ALEX. D\\
MATHEMATIQUES\qquad 5

\margin{154}
